\chapter{Materialism, Idealism, and Relational Ontology}

\section{The Strong Physicalist Claim}

Once descriptive adequacy has been separated from ontology, the materialism--idealism dispute can be stated more precisely. The strongest form of physicalism relevant here is not a doctrine of tiny solid objects. It is a sufficiency claim. Let $P_{\mathrm{phys},\le t}$ denote the complete public physical history required by the theory and let $E_t$ denote the experiential state. Strong physical supervenience asserts

\[
P_{\mathrm{phys},\le t}^{(1)}=P_{\mathrm{phys},\le t}^{(2)}
\quad\Longrightarrow\quad
E_t^{(1)}=E_t^{(2)}.
\]

No two systems can be physically identical in the complete relevant sense and experientially different. The corresponding causal-closure claim can be written schematically as

\[
p(Y_{t+\Delta}\mid P_{\mathrm{phys},\le t},E_t,do(A_t))
=
p(Y_{t+\Delta}\mid P_{\mathrm{phys},\le t},do(A_t)).
\]

Once the complete physical state is supplied, adding experience contributes no further information to physical prediction.

This formulation matters because it prevents the present mathematics from attacking a caricature. Relational states, complex waves, field descriptions, higher-order invariants, nonlocal mathematical representations, and dynamic symmetry breaking can all be accepted by a sophisticated physicalist. The physicalist need only maintain that these structures, whatever their final mathematical form, are the complete physical reality and are sufficient for experience. Wave ontology is therefore not idealism, and relational mathematics does not defeat physicalism merely by replacing object language with process language.

Naive substance materialism is more directly pressured. If persistent identity can be represented through relational organization even while components change, then immutable material bearers are not required for every form of identity. But this is a reorganization within possible physical ontology, not a proof that physical ontology fails.

\section{The Idealist Alternative}

Idealism assigns explanatory priority to mind, experience, intelligibility, or form. A relational idealist formulation can begin from a fundamental experiential-relational state $\Xi$ with dynamics

\[
\dot\Xi=G_I(\Xi),
\]

and treat the public physical world as a projection

\[
P_{\mathrm{phys},t}=\Pi(\Xi_{\le t}).
\]

The mathematical work in earlier chapters is compatible with such a view because invariant form can organize many changing realizations, but compatibility is not evidence. The fact that mathematical form plays a central role in classification does not establish that forms exist independently of physical realization. The fact that consciousness is epistemically immediate does not by itself determine the generative direction between experience and physical structure.

The Platonic material in this book should therefore be read carefully. Platonic symmetry offers an exact mathematical language of invariant form. It does not prove Platonic idealism. The affinity is philosophical: both emphasize structure that persists through changing appearance. The derivation itself remains neutral about whether the invariant is instantiated physically, mentally, or within some deeper ontology.

\section{Relational Dual-Aspect Realism}

A third possibility treats physical exteriority and experiential interiority as two potentially irreducible aspects of one deeper relational actuality. Let $r\in\mathscr R$ and define

\[
P_{\mathrm{phys}}=\Pi_{\mathrm{ext}}(r),
\qquad
E_{\mathrm{exp}}=\Pi_{\mathrm{int}}(r).
\]

Schematically,

\[
P_{\mathrm{phys}}\longleftarrow \mathscr R\longrightarrow E_{\mathrm{exp}}.
\]

This is not intended as substance dualism. The model contains one underlying relational state, not two independently interacting substances. Yet the diagram by itself is almost vacuous. Given sufficiently large latent space, one can always define a common container that stores two sets of variables. A scientifically meaningful dual-aspect model must therefore impose a nontrivial structural restriction on $\mathscr R$ and derive an observable consequence that cannot be reproduced equally well by simpler single-aspect models.

A local rank condition provides a minimal formal notion of complementarity. If $n=\dim\mathscr R$ and

\[
\operatorname{rank}D\Pi_{\mathrm{ext}}<n,
\qquad
\operatorname{rank}D\Pi_{\mathrm{int}}<n,
\]

neither projection locally reconstructs the underlying state. If the joint map satisfies

\[
\operatorname{rank}D(\Pi_{\mathrm{ext}},\Pi_{\mathrm{int}})=n,
\]

then the pair is locally sufficient even though each projection is individually incomplete.

This is a genuine mathematical notion of complementarity, but its semantic weakness must be stated immediately rather than after the machinery has acquired rhetorical momentum. Local rank complementarity establishes complementary coordinates or projections. It does not establish that one projection is phenomenal interiority. The interpretation of $E_{\mathrm{exp}}$ requires independent phenomenological anchoring.

\section{Dynamical Closure and Its Ambiguity}

Suppose the underlying state evolves by

\[
\dot r=F(r).
\]

The exterior projection obeys

\[
\dot P_{\mathrm{phys}}=D\Pi_{\mathrm{ext}}(r)F(r).
\]

For the projected physical state to possess a closed autonomous dynamics, there must exist $f$ such that

\[
\dot P_{\mathrm{phys}}=f(P_{\mathrm{phys}}).
\]

If two underlying states satisfy

\[
\Pi_{\mathrm{ext}}(r_1)=\Pi_{\mathrm{ext}}(r_2)
\]

but

\[
D\Pi_{\mathrm{ext}}(r_1)F(r_1)\neq D\Pi_{\mathrm{ext}}(r_2)F(r_2),
\]

then no exact autonomous law on that projection exists. A dual-aspect realization could in principle exhibit nonclosure of $P$ and nonclosure of $E$ separately while the pair $(P,E)$ closes.

The inference remains ambiguous because an incomplete physical coarse-graining produces the same mathematical symptom. Imagine a physical system with hidden phase $\theta$ and measured variable $x$. If $\dot x$ depends on $\theta$, then two states with the same $x$ but different hidden phase have different derivatives. The measured $x$ does not close, yet nothing nonphysical has been introduced. The failure simply tells us that the chosen physical description was incomplete.

This counterexample is central to the entire ontology audit. A failure of measured physical closure cannot be promoted directly into evidence against physicalism. The physical state must be enriched, alternative hidden physical variables must be tested, and approximate rather than exact closure must be considered because real experiments rarely recover every fast and slow variable exactly. The empirical program in Chapter 11 is designed around this problem.

\section{Difference Without Absolute Separation}

The relational framework also clarifies a philosophical confusion that often enters discussions of duality. Mathematical difference does not imply absolute ontological separation. In symbols, $A\neq B$ does not entail that $A$ and $B$ belong to unrelated substances or disconnected domains. Two projections may differ because they preserve different relational information. A state and its conjugate may differ while sharing density and symmetry. Two coordinates may differ while belonging to one underlying manifold. None of these examples proves dual-aspect monism, but all show that difference alone does not force substance dualism.

The ontological contest therefore remains genuinely open. Relational physicalism says the complete physical relational description is sufficient. Relational idealism assigns priority to experiential or intelligible relation. Relational dual-aspect realism proposes that one deeper actuality supports irreducible exterior and interior projections. The next chapters ask what evidence could discriminate among these claims without allowing the mathematics to decide by naming convention.
